\section{Connection elimination and constraint preservation}

We compare the connection-retaining and metric descriptions at the same perturbative order with the same action and endpoint data. The derivative-dependent LC substitution is not asserted to be a point canonical transformation on the unreduced space \cite{Pons2010}.

\stepcounter{equation}\begin{equation}\tag{D-01}\label{eq:D-01}
\frac{d\phi_A}{dt}=\{\phi_A,H_c\}+\{\phi_A,\phi_B\}\lambda^B\simeq0,\qquad H_T=H_c+\lambda^B\phi_B.
\end{equation}

In primary preservation, the image of the Poisson matrix fixes multipliers; only its left kernel produces new constraints. The original primary multipliers are retained at every consistency stage. Appendix C gives the scalar actions (D-02)--(D-03) and their closure.

\subsection{Direct comparison for a nonzero flat mode}

\begin{center}\small\begin{tabularx}{\textwidth}{@{}>{\raggedright\arraybackslash}X>{\raggedright\arraybackslash}X>{\raggedright\arraybackslash}X>{\raggedright\arraybackslash}X>{\raggedright\arraybackslash}X@{}}\toprule
Description & Configurations & First class & Second class & Pairs\\\midrule
Lorentz retained & 40 & 20 & 36 & 2\\
Lorentz fixed & 34 & 8 & 48 & 2\\
Metric & 10 & 8 & 0 & 2\\
\bottomrule\end{tabularx}\end{center}

The table is for one real Fourier quadrature. Retaining all 24 connection components and symbolic $k\ne0$, the reduced action, endpoint-derived two-form and the kinetic and gradient coefficients of both polarizations agree with (W-02) in all three descriptions. Including both real quadratures gives four real canonical pairs.

\stepcounter{equation}\begin{equation}\tag{D-05}\label{eq:D-05}
\Theta_{\rm ret}\big|_{\omega_{\rm LC}}-\Theta_{\rm met}=\delta B,
\qquad \Omega_{\rm ret}\big|_{\omega_{\rm LC}}-\Omega_{\rm met}=\delta^2B=0.
\end{equation}

$B$ includes the Lorentz total derivative and the complete NY endpoint. The current mixing the two real quadratures is retained, combining both NY terms from the same connection. For curved backgrounds we use the full Cartan solution, the auxiliary quadratic form and endpoint identities. We do not claim a new complete unreduced Dirac chain for every curved block; Appendix B specifies the directly verified scope.
